%%%%%%%%%%%%%%%%%%%%%%%%%%%%%%%%%%%%%%%%%%%%%%%%%%%%%%%%%%%%%%%%%%%%%
%%%%%%%%%%%%%%%%%%%%%%%%%%%%%%%%%%%%%%%%%%%%%%%%%%%%%%%%%%%%%%%%%%%%%
%%%%%%%%%%%%%%%%%%%%%%%%%%%%%%%%%%%%%%%%%%%%%%%%%%%%%%%%%%%%%%%%%%%%%
%%%%%%%%%%%%%%%%%%%%%%%%%%%%%%%%%%%%%%%%%%%%%%%%%%%%%%%%%%%%%%%%%%%%%
\chapter{Recording the conference}
%%%%%%%%%%%%%%%%%%%%%%%%%%%%%%%%%%%%%%%%%%%%%%%%%%%%%%%%%%%%%%%%%%%%%
%%%%%%%%%%%%%%%%%%%%%%%%%%%%%%%%%%%%%%%%%%%%%%%%%%%%%%%%%%%%%%%%%%%%%
%%%%%%%%%%%%%%%%%%%%%%%%%%%%%%%%%%%%%%%%%%%%%%%%%%%%%%%%%%%%%%%%%%%%%
%%%%%%%%%%%%%%%%%%%%%%%%%%%%%%%%%%%%%%%%%%%%%%%%%%%%%%%%%%%%%%%%%%%%%

%%%%%%%%%%%%%%%%%%%%%%%%%%%%%%%%%%%%%%%%%%%%%%%%%%%%%%%%%%%%%%%%%%%%%
%%%%%%%%%%%%%%%%%%%%%%%%%%%%%%%%%%%%%%%%%%%%%%%%%%%%%%%%%%%%%%%%%%%%%
\section{Recording vs. archiving}
%%%%%%%%%%%%%%%%%%%%%%%%%%%%%%%%%%%%%%%%%%%%%%%%%%%%%%%%%%%%%%%%%%%%%
%%%%%%%%%%%%%%%%%%%%%%%%%%%%%%%%%%%%%%%%%%%%%%%%%%%%%%%%%%%%%%%%%%%%%

At the beginning of this chapter, we distinguish between recording and archiving. The purpose of recording is to create a record that can later be reconstructed. Archiving additionally takes into account the overall context in which the record was created. We would like to emphasise this distinction in connection with the OAIS reference model \cite{OAIS-RM}, an ISO standard that provides a framework for long-term preservation. This rather extensive standard can be summarised by a simple, recursive idea:

\begin{quote}
Archiving data requires archiving the metadata and the tools needed to interpret that data. This process is repeated recursively, extending as far as the archival description of the protocols and encodings used.
\end{quote}

It is clear that archiving favours open technologies based on open standards over closed, undocumented systems. This leads us as far as the source code itself. Such an approach ensures that archived data can be reconstructed even several decades later.

A system based on freely available components is therefore the first step towards successful archiving. The remaining task is to address the recording process itself, define which aspects of communication should be archived, and select suitable tools.

%%%%%%%%%%%%%%%%%%%%%%%%%%%%%%%%%%%%%%%%%%%%%%%%%%%%%%%%%%%%%%%%%%%%%
%%%%%%%%%%%%%%%%%%%%%%%%%%%%%%%%%%%%%%%%%%%%%%%%%%%%%%%%%%%%%%%%%%%%%
\section{Data streams within an AccessGrid session}
%%%%%%%%%%%%%%%%%%%%%%%%%%%%%%%%%%%%%%%%%%%%%%%%%%%%%%%%%%%%%%%%%%%%%
%%%%%%%%%%%%%%%%%%%%%%%%%%%%%%%%%%%%%%%%%%%%%%%%%%%%%%%%%%%%%%%%%%%%%

AccessGrid consists of several interconnected components. Their communication can be divided into two main groups:
\begin{enumerate}
\item multimedia data transmission,
\item transmission of control messages.
\end{enumerate}

In Chapter \ref{kap:Multicast}, we described multicast as one of the delivery methods. It is particularly suitable for multimedia data in group communication. The RTP protocol is used for the actual transmission.

Control messages, on the other hand, account for only a small part of the overall communication and are therefore transmitted using the simpler unicast method. More specifically, communication between the VenueServer and the other AccessGrid components takes place in a form of XML-RPC.

\subsection{Candidates for archiving}

The most suitable candidates for recording, archiving and subsequent reconstruction are multimedia streams -- audio and video. The next candidates are the presentation display and, further, the desktop. Control messages and various message fragments related to security are less suitable candidates.

% In practice, it makes sense to record only streams that do not depend on state information. Successful recording would require state reconstruction, which may encounter limits of spatial and temporal complexity, and even computability in theoretical extremes.

Audio and video signals are essentially stateless. Codecs that generate transitions between keyframes rather than encoding individual frames introduce some complications. In the worst case, this can be handled by skipping to the next keyframe. The same principle applies to data fragmentation at individual ISO/OSI layers. The point is that audio and video carried in RTP packets should be relatively easy to archive. When reconstructing a damaged recording, the most complicated task will be to find the next keyframe, which depends on the codec being used.

Another candidate is the RFB protocol used for communication between a VNC server and client. The protocol is stateless, and terminating the connection or experiencing connection outages does not compromise the consistency of a VNC session.

As already mentioned, AccessGrid uses XML-RPC to control its components. In general, this communication is not of particular interest to us, with one exception: XML-RPC is used to synchronise AccessGrid shared applications. A shared presentation is an especially promising candidate for recording. One possible solution would be to store slide-transition events using a specially modified version of SharedPresentation. The problem, however, is the development of a universal recording tool for all shared applications. Moreover, the stored information could only be replayed in the same way in which it was created. In the context of OAIS, we would also have to archive all the viewers and shared applications used.

\subsection{Conversion to a common format}

\begin{quote}
\ldots or how to avoid $n^2$ conversion tools while using only $n$ formats
\end{quote}

To simplify the archiving process as much as possible, we decided to convert all events occurring during a videoconference into an audiovisual format and store them as RTP packets. Conversion may slightly reduce the quality of the original data, but it provides a uniform format for which playback and editing tools can later be developed.

\paragraph*{VNC display:} Conversion requires capturing the desktop image and encoding it with a suitable codec into RTP format.

\paragraph*{Shared presentation:} It is sufficient to capture the part of the desktop that displays the presentation. We can reuse the previous approach and use VNC to display the presentation. This solution is not limited to presentations in PDF or ODP format.

\begin{figure}
\centering
\includesvg[width=0.6\textwidth]{oais-environment}
\caption{The OAIS archival environment \cite{OAIS-RM}.}
\label{fig:OaisEnvironment}
\end{figure}

The OAIS reference model defines four elements of an archival environment (Figure \ref{fig:OaisEnvironment}). From our perspective, the \textquotedblleft Producer\textquotedblright\ is any application that transmits RTP packets, namely VIC and RAT. The \textquotedblleft Consumer\textquotedblright\ is the user who wants to reconstruct and play back an archived recording, for example using VIC and RAT. To follow the reference model, we would also have to provide a \textquotedblleft Management\textquotedblright\ entity responsible for setting the policies followed by the OAIS archival environment. Archive maintenance is therefore not performed by management itself, but by structures within the archive.

Implementing the OAIS model is outside the scope of this thesis, and the remainder of this chapter is independent of the archival paradigm used.

%%%%%%%%%%%%%%%%%%%%%%%%%%%%%%%%%%%%%%%%%%%%%%%%%%%%%%%%%%%%%%%%%%%%%
%%%%%%%%%%%%%%%%%%%%%%%%%%%%%%%%%%%%%%%%%%%%%%%%%%%%%%%%%%%%%%%%%%%%%
\section{Existing recording tools}
%%%%%%%%%%%%%%%%%%%%%%%%%%%%%%%%%%%%%%%%%%%%%%%%%%%%%%%%%%%%%%%%%%%%%
%%%%%%%%%%%%%%%%%%%%%%%%%%%%%%%%%%%%%%%%%%%%%%%%%%%%%%%%%%%%%%%%%%%%%

\subsection{Tcpdump -- packet sniffer}
Tcpdump can save packets passing through a network interface. At first glance, it might seem to be an ideal tool for capturing all information in a data transmission. The problem with recording RTP packets, however, is preserving their timing. Although the resulting file contains all audiovisual data, reconstructing the original timing is difficult. Tcpdump is therefore not a suitable tool for this purpose.

\subsection{RTPtools}
RTPtools is a set of utilities for working with RTP streams. It includes rtprecord, rtpplay and rtpproxy. Another tool is the RTP reflector, which can mirror incoming data, mix different RTP streams and transparently transcode communication using another codec.
All these tools are command-line applications and provide only basic functionality.

\subsection{AGVCR -- RTP recorder/player}
At the time of writing, the most suitable tool for recording RTP packets is AGVCR by D. Piper \cite{AGVCR-home} of Indiana University. It provides both graphical and text-based interfaces (including Telnet), and supports automatic launching of applications used for recording and playing multimedia. The recording is stored in a separate file containing information about all received data streams and their timing. The tool also allows a subset of the recorded data to be exported by splitting it into smaller segments.

%%%%%%%%%%%%%%%%%%%%%%%%%%%%%%%%%%%%%%%%%%%%%%%%%%%%%%%%%%%%%%%%%%%%%
%%%%%%%%%%%%%%%%%%%%%%%%%%%%%%%%%%%%%%%%%%%%%%%%%%%%%%%%%%%%%%%%%%%%%
\section{Recording an ongoing conference}
%%%%%%%%%%%%%%%%%%%%%%%%%%%%%%%%%%%%%%%%%%%%%%%%%%%%%%%%%%%%%%%%%%%%%
%%%%%%%%%%%%%%%%%%%%%%%%%%%%%%%%%%%%%%%%%%%%%%%%%%%%%%%%%%%%%%%%%%%%%

The purpose of recording is to capture the course of a conference as faithfully as possible from the perspective of a single node. The simplest option, which we also used for the node at the Ciant organisation, is to record the signal from a graphics card, for example as an S-Video signal. The advantages and disadvantages of this solution were discussed in Chapter \ref{kap:CiantSkusenosti}.

A more complex approach uses AGVCR. Recording requires a dedicated computer because, within one operating system, it is not possible to record, view and transmit packets using the same socket at the same time. The last application to call \verb|bind()| receives the packets; the other applications can only transmit data using \verb|send()| or \verb|sendto()|.

For a large conference, the recording node must have sufficient processing power, disk space and a high-quality network connection. It stores all packets appearing in the conference. Fortunately, the node does not need to encode or decode any data streams; its only task is to store RTP packets.

%%%%%%%%%%%%%%%%%%%%%%%%%%%%%%%%%%%%%%%%%%%%%%%%%%%%%%%%%%%%%%%%%%%%%
%%%%%%%%%%%%%%%%%%%%%%%%%%%%%%%%%%%%%%%%%%%%%%%%%%%%%%%%%%%%%%%%%%%%%
\section{StreamVNC -- a new tool for local recording}
%%%%%%%%%%%%%%%%%%%%%%%%%%%%%%%%%%%%%%%%%%%%%%%%%%%%%%%%%%%%%%%%%%%%%
%%%%%%%%%%%%%%%%%%%%%%%%%%%%%%%%%%%%%%%%%%%%%%%%%%%%%%%%%%%%%%%%%%%%%

The problem of conference recording led to the idea of creating a tool that would allow a presentation to be prepared by transmitting it to a local network interface -- localhost. During the design of the system, we considered the following requirements:
\begin{paklist}
\item Simple operation.
\item A graphical user interface.
\item The ability to record the desktop, for example while a presentation is running.
\item Sufficient recording quality.
\item Low network bandwidth requirements.
\item Support for both audio and video recording.
\item Reuse of existing components.
\item Easy editing (cutting) of the resulting recording.
\item Simple playback.
\end{paklist}

The first problem was finding suitable components that could communicate with one another and were sufficiently configurable and controllable. They had to cover the following areas:
\begin{paklist}
\item Recording RTP packets.
\item Audio processing.
\item Capturing and processing camera video.
\item Capturing and processing the desktop image.
\item Providing suitable presentation tools.
\end{paklist}

The model situation is shown in Figure \ref{fig:LocalhostRecording}. The key element is the VNC system (1), which can display an application (2) independently of the surrounding desktop. The figure shows the VNC content being sent to a local address by the VIC application (3), together with audio from RAT (4) and an additional camera video stream (5). All three data streams are recorded by AGVCR (6).

\begin{figure}
\centering
\includesvg[width=0.8\textwidth]{localhost-recording}
\caption{Preparing a presentation for later use.}
\label{fig:LocalhostRecording}
\end{figure}

\subsection{Recording RTP packets}
The excellent AGVCR tool simplified the selection of a suitable recording application. It required no further modifications because it was highly stable, efficient and configurable. The only missing feature was the ability to cut a segment from a saved recording. At the time of writing, we had already submitted this feature request to the application author.

\subsection{Audio processing}\label{kap:HybridRAT}
Before use, we had to make minor modifications to the RAT application, which transmits audio using RTP packets. The standard RAT 4.3 version from UCL had a problem when a firewall was enabled. The application refused to start, which was unacceptable for our purposes. The cause was communication by the \textquotedblleft mbus\textquotedblright\ component used by RAT. At startup, it always attempted to contact one of the multicast addresses, causing the application to freeze when a firewall was running. After some time, we found another modified version of RAT from VRVS, which started without problems.

Its disadvantage was the lack of an interface for ALSA (Advanced Linux Sound Architecture). It supported only OSS (Open Sound System), which had poor performance and problems interacting with other applications using audio. We therefore used the source code of both implementations to create a hybrid version, RAT-SVNC, combining the advantages of each.

\subsection{Desktop video processing}
Within AccessGrid, the VIC application is used to process video signals. Again, we could not use the original version from UCL. The application locks the socket on which it transmits RTP packets, preventing AGVCR from recording them. It also includes a poorly functioning H.263 codec, whose use often causes a \textquotedblleft segmentation fault\textquotedblright\ error. The VRVS version of VIC, based on the OpenMash project, contained a working H.263 codec and had a more compact design, which is particularly important for end users. In the EVO project, we also found a newer version of VIC called Vievo. We liked its design and decided to use it as the basis for the new tool.

We chose the name EVIC for the new version of VIC. It contains the core from the OpenMash-based VIC-VRVS and parts of the Vievo design. We added controls that allow recording to be controlled from an external X11 display. We use this feature when recording VNC content, allowing the controls to be placed on the physical X11 display.

The performance of the OpenMash VIC allows us to achieve 2 fps at PAL resolution and 10 fps at CIF resolution. These values are sufficient for presentation purposes and also save network bandwidth. Nevertheless, the performance of this component should be improved in the future. Tests showed that 8 fps at PAL resolution used 100\% of the processing capacity of a 1.8 GHz Pentium 4 processor.

\subsection{Camera video processing}

We also successfully used VIC to capture camera video at 25 fps and CIF resolution. PAL resolution is not yet supported. Camera video is intended as supplementary information accompanying the more important presentation and audio.

\subsection{Support for presentation tools}

When selecting presentation tools, we focused primarily on support for PDF, OpenDocument Format (ODF) and PowerPoint (PPT). We also wanted an image viewer, web-page presentation, limited support for video sequences and support for the VMware virtual machine. Extensibility to additional applications was also essential. The version of TightVNC we use does not support 3D acceleration, so applications using OpenGL are not supported yet.

In addition to the application itself, the VNC display required a suitable window manager. Important selection criteria included its size, resource requirements and the placement of window decorations. We eventually selected FLWM, which places window decorations on the left side. This avoids visual conflicts with the decorations of the VNC viewer. Before using it, we made minor modifications to FLWM, including adding the \verb|-e| parameter, which allows another program to be started when FLWM launches. For example, this allows us to set the desktop background using the command \verb|xsetroot -solid '#004455'|.

\subsection{StreamVNC architecture}

All components mentioned above were integrated into a single system, StreamVNC, shown in Figure \ref{fig:ArchStreamVNC}.

The system offers a high degree of configurability while requiring very little from the user. It uses C, TCL/Tk and shell script. It requires a running X server, on which all controls are displayed. The main architectural components are:
\begin{paklist}
\item RTP packet recording,
\item audio capture,
\item VNC -- server, viewer and desktop capture,
\item camera video capture,
\item applications.
\end{paklist}

\begin{figure}
\centering
\includesvg[width=1\textwidth]{arch-stream-vnc}
\caption{StreamVNC system architecture.}
\label{fig:ArchStreamVNC}
\end{figure}

Using Xvnc (the TightVNC server) requires the \verb|fixed| and \verb|helvetica| fonts. The font path is usually configured correctly by the Linux distribution, as is the case with OpenSUSE. Some distributions, such as Gentoo, require the font path to be set explicitly using \verb|Xvnc -fp <dir>|. StreamVNC therefore provides a binary version of Xvnc together with the essential fonts and the \verb|fonts.dir| and \verb|fonts.alias| definitions. The Xvnc server is not started directly, but through the \verb|vnc/Xvnc-run| script, which sets the custom font path. If necessary, this script can simply be modified.

\subsection{Implemented profiles}
The main script, \verb|agvcr-run|, is started either without parameters or with one parameter specifying the name of the file to be played. The application therefore distinguishes two profiles using separate configuration files: recording uses \verb|.agvcr.conf|, while playback uses \verb|.play.conf|.

\paragraph*{Recording profile:} This is the default profile, intended for recording a presentation. Programs shown in grey in Figure \ref{fig:SequenceStreamVNC} are started only in this profile. The user is presented with a menu from which one of the applications can be selected.

\paragraph*{Playback profile:}
In this profile, AGVCR automatically transmits RTP packets from the file supplied as its parameter. Only one instance of EVIC is started, after it receives an RTCP packet. The VNC server and VNC viewer are not started because they would unnecessarily occupy screen space during playback. The transmitted RTP packets are automatically received by RAT and EVIC, and the user sees the content of the saved presentation.

\begin{figure}
\centering
\includesvg[width=1\textwidth]{sequence-stream-vnc}
\caption{Dependencies between programs started by StreamVNC.}
\label{fig:SequenceStreamVNC}
\end{figure}

\subsection{Managing running processes}
Since this is an integrated system consisting of many communicating processes, process management was one of the most important issues. The main executable, called \verb|stream-vnc|, is implemented as a shell script. Its task is to start child processes, monitor their activity and terminate them all when necessary. The main obstacle was the limited flexibility of shell job control\footnote{Specifically, GNU bash version 3.00.16(1)-release (i586-suse-linux).}. We eventually solved the problem as follows:

The main program maintains a list of processes, \verb|$REQ|, whose termination should cause the entire system to shut down. It registers a signal handler for \verb|SIGCHLD|. When a child process terminates, the operating system kernel automatically notifies the parent process (signal 17). The shell uses the \verb|jobs -pr| command to obtain the list of running child processes. If a process from the \verb|$REQ| list is missing, the shutdown sequence is started and the parent sends \verb|SIGTERM| to the remaining child processes. In addition, all processes using the VNC display are automatically terminated when the VNC server exits.

Another approach we attempted to implement used the terminal functionality provided by the operating system. A terminal is a special device intended for inter-process communication. The operating system provides system calls for creating and managing this device. The process that creates the device becomes the leader of the process group that communicates through it. The group is identified by its process-group ID, which is the PID of its leader. If the leader terminates, the operating system should send the \verb|SIGHUP| signal to all processes in the group. The advantage is automatic signal delivery, unlike our application-level solution in the shell script.

The disadvantage is that some versions of shell interpreters do not provide equivalents of the \verb|grantpt()|, \verb|unlockpt()| and \verb|ptsname()| system calls. This must be handled by a small C program. It opens the special file \verb|/dev/ptmx|, which creates the master side of the communication channel. The Unix API provides the functions \verb|posix_openpt()|, \verb|grantpt(master)|, \verb|unlockpt(master)| and \verb|ptsname(master)|. The last function returns the name of a file such as \verb|/dev/pts/X|, where \verb|X| represents the slave side of the communication channel. A simple program therefore opens a pseudo-terminal, sets its file descriptors correctly, calls \verb|setsid()| to create a new session and process group, and becomes the leader of that group. It can then call one of the \verb|exec()| functions to start another process in its address space. Practical experience with StreamVNC showed, however, that the mechanism we selected is sufficient.

\subsection{Environment consistency check}
To make it easier to identify problems caused by missing applications, we added the \verb|--check| parameter to the \verb|stream-vnc| script. Before the first use, it is therefore advisable to run \verb|./stream-vnc --check|. This checks the consistency of the system and identifies the most common problems. Among other things, it also detects whether the operating system kernel's \verb|vmmon| module for VMware support is running.

\proofed

\begin{comment}

%%%%%%%%%%%%%%%%%%%%%%%%%%%%%%%%%%%%%%%%%%%%%%%%%%%%%%%%%%%%%%%%%%%%%
%%%%%%%%%%%%%%%%%%%%%%%%%%%%%%%%%%%%%%%%%%%%%%%%%%%%%%%%%%%%%%%%%%%%%
\section{Processing archived content}
%%%%%%%%%%%%%%%%%%%%%%%%%%%%%%%%%%%%%%%%%%%%%%%%%%%%%%%%%%%%%%%%%%%%%
%%%%%%%%%%%%%%%%%%%%%%%%%%%%%%%%%%%%%%%%%%%%%%%%%%%%%%%%%%%%%%%%%%%%%

How to overlay an image on a video:
http://www.videolan.org/doc/play-howto/en/ch04.html
\begin{verbatim}
vlc test.ogm --sub-filter logo --logo-file logo.png --logo-x 20 --logo-y 20
\end{verbatim}

Real-time transcoding to MPEG-4 video (128 kb/s) and OGG audio (32 kb/s):
\begin{verbatim}
vlc -vvv test.ogg --sout '#transcode{vcodec=mp4v, acodec=vorbis,vb=128,ab=32,deinterlace}:display'
\end{verbatim}

Various filters are described at:
http://www.videolan.org/doc/streaming-howto/en/ch03.html

2xPIP

\begin{verbatim}
vlc -vvv test.ogg --sout '#duplicate{
dst={transcode{vcodec=mp2v,vb=4096,scale=1,sfilter=mosaic}:bridge-in{offset=100}:display},
dst={duplicate{dst=mosaic-bridge{id=1,width=320,height=200},dst=mosaic-bridge{id=2,width=320,height=200}}}}'
\end{verbatim}

\begin{verbatim}
vlc --color -I telnet --vlm-conf /dev/shm/mosaic.conf --fake-file background.jpg
--fake-aspect-ratio "4:3" --mosaic-width 640 --mosaic-height 400 --mosaic-keep-picture --mosaic-rows 2 --mosaic-cols 2 --mosaic-position 1 --mosaic-order 1,2,3

MPlayer V4L loopback device
http://lists.mplayerhq.hu/pipermail/mplayer-dev-eng/
001-November/002600.html

vloopback
http://www.lavrsen.dk/twiki/bin/view/Motion/
VideoFourLinuxLoopbackDevice

VJ software that supports vloopback input/output:
http://www.gephex.org/

##################################
## HTTP interface mosaic wizard ##
##################################

# Comment out the following line if you do not want to reset the VLM configuration
del all

# Background options
new   bg broadcast enabled
setup bg input     "fake:" option "fake-file=background.jpg" option "fake-width=640" option "fake-height=400" option "fake-ar=4:3"
setup bg output\
	#transcode{vcodec=mp2v,vb=4096,scale=1,sfilter=mosaic}:\
	bridge-in{offset=100}:

# Mosaic options
setup bg option mosaic-alpha=200
setup bg option mosaic-height=400
setup bg option mosaic-width=640
setup bg option mosaic-align=5
setup bg option mosaic-xoffset=10
setup bg option mosaic-yoffset=10
setup bg option mosaic-vborder=5
setup bg option mosaic-hborder=10
setup bg option mosaic-position=1
setup bg option mosaic-rows=2
setup bg option mosaic-cols=2
setup bg option mosaic-order=1,2,3
setup bg option mosaic-delay=0
setup bg option mosaic-keep-picture

new channel1 broadcast enabled
setup channel1 input "test.ogg"
setup channel1 output \
	#mosaic-bridge{id=1,width=320,height=200},\
	select=audio,\
	dst=bridge-out{id=3},\
	select=audio

new channel2 broadcast enabled
setup channel2 input "test.ogg"
setup channel2 output \
	#duplicate{
		dst=mosaic-bridge{id=2,width=320,height=200},\
		select=video,\
		dst=bridge-out{id=3},\
		select=audio\
	}

# Input options
# Launch everything

control bg play
control channel1 play
control channel2 play

# end of mosaic batch
\end{verbatim}

RTSP is a text-based protocol similar to HTTP that could be used for media streaming. LIVE555 is a library that could provide this functionality.
Presentations could be stored in an appropriate format and played on demand by a server. Several solutions already exist, such as FFmpeg, Helix DNA and VLC.

An unresolved problem is transcoding RTP streams stored in the rtpplay format. The AGVCR format has the same problem. Ideally, we would need a utility that separates these streams into individual files according to their SSRC. We would then need another utility to convert each file into a suitable format that could be processed by mencoder or FFmpeg. VLC could then be used to create a grid containing several such streams.

The VRVS semester project indicates that XviD is suitable for archiving, but that no method was found for playing back a recorded RTP stream.

The thesis DP2006Patoraj.pdf describes JRAT, which is faster than RAT. It is written in Java but consumes somewhat more memory.

The thesis DP2006Svrcek.pdf describes an attempt to modify VIC so that it stores the stream in a file in YUV format, which is then compressed to XviD.

The following configuration was intended to record directly from VIC using MPlayer, but it did not work correctly. In theory, the procedure should be as follows:
Create an SDP file named test.sdp:
\begin{verbatim}
c=IN IP4 127.0.0.1
m=video 10000 RTP/AVP 31
\end{verbatim}

Then run mencoder with \verb|-ovc [raw|lavc|copy]| and \verb|sdp://test.sdp|, or try MPlayer instead of mencoder.

Next, start \verb|vic 127.0.0.1/10000|. The video appears in MPlayer, but with artefacts, and the log contains errors such as:
\begin{verbatim}
Error while decoding frame!
[h261 @ 0x8861b78]Bad picture start code
[h261 @ 0x8861b78]header damaged
Error while decoding frame!
[h261 @ 0x8861b78]Bad picture start code
[h261 @ 0x8861b78]header damaged
Error while decoding frame!
[h261 @ 0x8861b78]run overflow at 15x5
[h261 @ 0x8861b78]Error at MB: 130
\end{verbatim}

Conclusion: this approach does not work. VLC and FFmpeg exhibit the same problem.

\end{comment}
